\chapter{Immunisation During Pregnancy}

\section*{Definition and Overview}
Immunisation during pregnancy protects both the mother and her baby from serious infections. Two vaccines are recommended for every pregnancy in Australia: the whooping cough (pertussis) vaccine, given in the third trimester, and the influenza (flu) vaccine, given at any stage of pregnancy. COVID-19 vaccination is also recommended. These vaccines work in a unique way during pregnancy: the mother produces antibodies that cross the placenta and provide the newborn with passive protection during the first months of life, when babies are too young to be vaccinated themselves and are most vulnerable to severe disease. Vaccination in pregnancy is safe and well studied, and it is one of the most effective measures available for protecting infants.

\section*{Epidemiology}
Influenza and whooping cough are significant causes of illness in pregnant women and newborns. Influenza increases the risk of severe respiratory illness, hospitalisation, and complications in pregnancy, and it can be life-threatening to both mother and baby. Whooping cough is particularly dangerous for young infants: most hospitalisations and deaths from the disease in Australia occur in babies under six months of age. Rates of vaccination in pregnancy have improved over time, and in recent years more than three-quarters of Australian women have received the whooping cough and influenza vaccines during pregnancy, contributing to a marked reduction in infant whooping cough. COVID-19 vaccination in pregnancy has also become an established recommendation.

\section*{Aetiology and Pathophysiology}
During pregnancy, maternal antibodies are actively transported across the placenta, particularly in the third trimester, providing the newborn with the same antibodies the mother has produced. This is the key to vaccination in pregnancy: when a woman is vaccinated, she develops antibodies that both protect her from the infection and are passed on to her baby. Antibodies against influenza and whooping cough are transferred efficiently, and high maternal antibody levels in the weeks before birth ensure the baby has protective levels at delivery and for the first two to three months of life. The whooping cough vaccine is timed for 20--32 weeks of pregnancy to maximise antibody transfer, and re-vaccination is recommended in every pregnancy because antibody levels fall between pregnancies.

\section*{Clinical Features}
\begin{itemize}
    \item Recommended vaccines in pregnancy: pertussis (whooping cough), influenza, and COVID-19
    \item Whooping cough vaccine: given between 20 and 32 weeks of pregnancy, ideally in the early part of this window
    \item Influenza vaccine: recommended at any stage of pregnancy, before or during the flu season
    \item Vaccines are free for pregnant women under the National Immunisation Program
    \item Other vaccines may be recommended for travel or medical reasons, with timing discussed with a doctor
    \item Mild local reactions such as soreness at the injection site are common; the vaccines do not contain live virus and cannot infect the baby
\end{itemize}

\section*{Differential Diagnosis}
Some vaccines are not used in pregnancy, and individual risk assessment is needed.
\begin{itemize}
    \item \textbf{Live vaccines:} measles-mumps-rubella (MMR) and varicella vaccines are avoided in pregnancy because of theoretical risk to the fetus
    \item \textbf{Vaccines for specific risks:} hepatitis B, hepatitis A, and meningococcal vaccines can be given in pregnancy when there is a clear indication
    \item \textbf{Travel vaccines:} timing depends on destination, risk, and the specific vaccine
    \item \textbf{Catch-up vaccines:} some vaccines can be given after birth to complete schedules
    \item \textbf{Household and carer vaccination:} vaccinating close contacts, including the father, reduces infant exposure ("cocooning")
\end{itemize}

\section*{When to Seek Medical Care}
Pregnant women should discuss vaccination with their GP, midwife, or obstetrician at an early antenatal visit and plan the timing of recommended vaccines. Ask about catch-up vaccination if a vaccine was missed in a previous pregnancy or if the schedule is incomplete. Any severe or concerning reaction to a vaccine should be reported, although serious reactions are extremely rare.

\textbf{Call 000 (or local emergency services) for:}
\begin{itemize}
    \item Difficulty breathing, wheezing, or swelling of the face and lips within minutes to hours of a vaccine (anaphylaxis, which is very rare)
    \item Fainting with injury
    \item Any severe allergic reaction symptoms
\end{itemize}

\section*{Diagnosis and Investigations}
\begin{itemize}
    \item \textbf{Antenatal record review:} checking immunisation history and identifying what is due in pregnancy
    \item \textbf{Serology:} blood tests for immunity in specific situations, such as rubella and varicella screening in early pregnancy
    \item \textbf{Documentation:} vaccination recorded on the Australian Immunisation Register
    \item \textbf{Risk assessment:} review of travel, occupation, and medical conditions that may warrant additional vaccines
\end{itemize}

\section*{Management}
\begin{itemize}
    \item \textbf{Pertussis vaccine:} one dose in every pregnancy, given at 20--32 weeks; ideally early in this window to maximise antibody transfer
    \item \textbf{Influenza vaccine:} one dose per pregnancy, inactivated vaccine, at any trimester before or during the flu season
    \item \textbf{COVID-19 vaccine:} recommended at any stage of pregnancy, with boosters as advised
    \item \textbf{Timing for other vaccines:} planned individually based on risk and vaccine type
    \item \textbf{Postnatal catch-up:} women who miss vaccination in pregnancy can receive the vaccines after birth, protecting themselves and helping to protect their baby
    \item \textbf{Partner and family vaccination:} encouraged to protect the newborn
\end{itemize}

\section*{Complications}
Vaccination in pregnancy is associated with a very small number of complications, and the risk of disease far outweighs vaccine risk.
\begin{itemize}
    \item Common, mild reactions: soreness, redness, or swelling at the injection site, low-grade fever, and fatigue
    \item Rare allergic reactions, including anaphylaxis
    \item Unvaccinated women and their babies remain at risk of influenza, whooping cough, and their complications
    \item Influenza in pregnancy increases the risk of pneumonia, hospitalisation, and preterm birth
    \item Whooping cough in a newborn can cause severe breathing difficulty, pneumonia, and death
\end{itemize}

\section*{Prognosis and Outcomes}
Vaccination in pregnancy is highly effective. Maternal influenza vaccination substantially reduces the risk of the mother developing influenza and reduces the risk of influenza and hospitalisation in the infant during the first six months of life. Whooping cough vaccination in pregnancy reduces the risk of whooping cough in young infants by more than 90\%, and infant deaths from the disease have fallen markedly since the program was introduced. Antibodies transferred in pregnancy protect the baby until routine childhood vaccination begins at six to eight weeks of age. These are among the best-evidenced preventive measures in maternity care.

\section*{Prevention and Self-Care}
\begin{itemize}
    \item Book the whooping cough vaccine for 20--32 weeks of pregnancy as soon as the pregnancy is confirmed
    \item Have the influenza vaccine at any stage of pregnancy, ideally before winter
    \item Ask your GP or midwife about COVID-19 vaccination and boosters
    \item Ensure the household and close family members are up to date with their vaccinations
    \item Keep a record of vaccines given in pregnancy
    \item Remember that vaccines given in pregnancy are safe for both mother and baby
\end{itemize}

\section*{Sources and Further Reading}
The following reputable sources provide further information for healthcare students and patients seeking reliable, current advice about immunisation during pregnancy.
\begin{itemize}
    \item Australian Government Department of Health and Aged Care. \textit{Immunisation during pregnancy.} https://www.health.gov.au/
    \item Healthdirect Australia. \textit{Vaccination in pregnancy.} https://www.healthdirect.gov.au/
    \item Pregnancy, Birth and Baby. \textit{Vaccinations during pregnancy.} https://www.pregnancybirthbaby.org.au/
    \item Royal Australian and New Zealand College of Obstetricians and Gynaecologists (RANZCOG). \textit{Vaccination in pregnancy.} https://ranzcog.edu.au/
    \item NHS. \textit{Vaccinations in pregnancy.} https://www.nhs.uk/pregnancy/
    \item Centers for Disease Control and Prevention. \textit{Vaccines during pregnancy.} https://www.cdc.gov/
\end{itemize}
